% SPDX-License-Identifier: Apache-2.0
% covers: Dmem
\datasheet{Dmem}{Vreteno's data memory: 1024 words behind an AXI peripheral end.}

\dsfacts{\code{cpu/vreteno/src/dmem.rs}}{\code{vreteno32::dmem::Dmem}}%
{\code{//docs:vreteno}}

\subsection*{Function}

\code{Dmem} is the data memory of the Vreteno machine, a device on its
bus. It holds \code{DMEM\_WORDS}, 1024 words, in four memories of a
byte, \code{lane0} to \code{lane3}, one per lane of the word. A store
of a byte or a half writes only its lanes and reads nothing.

It is the peripheral client of an \code{AxiPer}, and its link is one
port, \code{bus}, a \code{PerPort} of the four channels: it takes
requests on \code{bus\_req} and write beats on \code{bus\_w}, and
answers on \code{bus\_ans} for a write and \code{bus\_r} for a read.
It answers single-beat bursts,
which is all the core makes.

\subsection*{Parameters}

\begin{center}\footnotesize
\begin{tabular}{@{}lp{0.7\textwidth}@{}}
\toprule
Parameter & Meaning \\
\midrule
\code{I} & The width of the AXI identifier on \code{bus\_req},
\code{bus\_ans} and \code{bus\_r}, in bits. The tables use 4, as the board
does since its bus has two hosts behind an arbiter (issue 241). \\
\bottomrule
\end{tabular}
\end{center}

\subsection*{Address map}

The word a burst names is bits 11 to 2 of its address. Lane $k$ is
byte $k$ of the word: \code{lane0} is bits 7 to 0, \code{lane3} bits 31
to 24. The memory looks at no other address bits.

On the board, \code{BoardRouter} sends it the page at \code{0x1000}
with mask \code{0xffff\_f000}, so addresses \code{0x1000} to
\code{0x1ffc} are words 0 to 1023. The model's \code{DATA\_BASE} is
the same \code{0x1000}.

\dstables{Dmem}

\subsection*{Behaviour}

\begin{itemize}
\item A write is taken when no other write is waiting for its beat.
Its word and identifier go into \code{paddr} and \code{pid}, and
\code{pend} is set.
\item The held write's beat is taken when it has arrived and
\code{bus\_ans} has room. Each lane whose strobe bit is set takes its byte
of the beat, and \code{bus\_ans} sends \code{OKAY} with \code{pid}.
\item A read is taken when no write is waiting and the answer register
is free or is being emptied in this cycle. The four lanes are read
into \code{word} at the edge, the synchronous read a block RAM has.
\item The read is answered on \code{bus\_r} the cycle after, with
\code{rid}, \code{OKAY} and \code{last} set.
\item While a write waits for its beat, no request of either kind is
taken.
\item \code{Dmem::with} fills the lanes from bytes before the first
cycle. \code{data\_word} reads a word back for a run or a test.
\end{itemize}

\subsection*{Verification}

\begin{itemize}
\item \code{//cpu/vreteno:lockstep\_test} compares every word of the model's
data memory with the four lanes at the halt, in \code{demo\_program},
\code{a\_multiply\_right\_after\_a\_load} and \code{random\_programs}.
\item \code{//cpu/vreteno:hello\_test}, \code{//cpu/vreteno:hello\_cc\_test}
and \code{//cpu/vreteno:board\_test} put the compiled programs'
constants in it with \code{Dmem::with}, and check what the programs
print.
\item The build simulates \code{Dmem<2>}'s netlist against the trace
of \code{//cpu/vreteno:vreteno} under nvc and Verilator:
\code{//docs:vreteno\_sim\_dmem}\allowbreak\code{\_dmem\_tb\_test} and
\code{//docs:vreteno\_vsim\_dmem\_test}.
\end{itemize}

\subsection*{Used by}

\code{Board}, as \code{dmem} behind the tracker \code{pdmem}. The
demonstration \code{//cpu/vreteno:vreteno}, the lockstep test and
\code{run.rs}. \code{board\_netlist} writes a program's data into
\code{lane0} to \code{lane3} of the board's netlist.

\subsection*{Limits}

It serves bursts of one beat and does not look at a burst's length. It
checks no address, since the router sends it only its own range, so
an address beyond word 1023 wraps. One write is held at a time. Nothing
on the machine fills it at run time except stores, so a program's
constants have to be in it before the first cycle.
\code{cpu/vreteno/board/ax7a200.xdc} constrains the four lanes to block
RAM.
